\subsection{形式幂级数系统的合成}

设 A 为一个交换环；称一个系统

\[
\mathbf {f} = (f _ {1}, \dots , f _ {p}) \in (\mathrm{A} [ [ \mathrm{X} _ {1}, \dots , \mathrm{X} _ {q} ] ]) ^ {p}\tag{8}
\]

在 \(X_{j}\) 中（ \(1 \leqslant j \leqslant q\) ）、系数在 A 中的形式幂级数系统，如果所有 \(f_{j}\) 都无常数项，则称其无常数项。对于每个形式幂级数系统 (8) 以及每个系统

\[
\mathbf {g} = (g _ {1}, \dots , g _ {q}) \in (\mathrm{A} [ [ \mathrm{X} _ {1}, \dots , \mathrm{X} _ {r} ] ]) ^ {q}\tag{9}
\]

由 \(q\) 个无常数项的形式幂级数组成的系统，我们将 \(\mathbf{f} \circ \mathbf{g}\)（或 \(\mathbf{f}(\mathbf{g})\)）记作由 \(f_{j}(g_{1},\ldots ,g_{q})\)（\(1\leqslant j\leqslant p\)）组成的、属于

\[
\left. \left(\mathrm{A} \left[ \left[ \mathrm{X} _ {1}, \dots , \mathrm{X} _ {r} \right] \right]\right) ^ {p} \right.
\]

（《代数》，第 IV 章，§ 5，第 5 节）。如果

\[
\mathbf {h} = (h _ {1}, \dots , h _ {r}) \in (\mathrm{A} [ [ \mathrm{X} _ {1}, \dots , \mathrm{X} _ {s} ] ]) ^ {r}
\]

是第三个无常数项的系统，则

\[
(\mathbf {f} \circ \mathbf {g}) \circ \mathbf {h} = \mathbf {f} \circ (\mathbf {g} \circ \mathbf {h}).\tag{10}
\]

对于每个整数 m，

\[
\left(\mathbf {f} ^ {(m)} \circ \mathbf {g} ^ {(m)}\right) \circ \mathbf {h} ^ {(m)} = \mathbf {f} ^ {(m)} \circ \left(\mathbf {g} ^ {(m)} \circ \mathbf {h} ^ {(m)}\right)
\]

其中，\(\mathbf{f}^{(m)}\)、\(\mathbf{g}^{(m)}\)、\(\mathbf{h}^{(m)}\) 表示形式幂级数系统 f、g、h 中由总次数 \(\leqslant m\) 的项组成的多项式系统。但显然，级数中总次数 \(\leqslant m\) 的项，在 \((\mathbf{f} \circ \mathbf{g}) \circ \mathbf{h}\)（分别在 \(\mathbf{f} \circ (\mathbf{g} \circ \mathbf{h})\)）中，与 \((\mathbf{f}^{(m)} \circ \mathbf{g}^{(m)}) \circ \mathbf{h}^{(m)}\)（分别与 \(\mathbf{f}^{(m)} \circ (\mathbf{g}^{(m)} \circ \mathbf{h}^{(m)})\)）相同，故断言成立。

对于每个系统 (8)，我们以 \(M_{\mathbf{f}}\) 或 \(M_{\mathbf{f}}(\mathbf{X})\) 表示雅可比矩阵 \((\partial f_i / \partial \mathbf{X}_j)\)（ \(1 \leqslant i \leqslant p, 1 \leqslant j \leqslant q\)），其中 \(i\) 是行指标，\(j\) 是列指标；对于两个系统 (8) 和 (9)，若 \(\mathbf{g}\) 无常数项，则

\[
M _ {\mathbf {f} \circ \mathbf {g}} = (M _ {\mathbf {f}} (\mathbf {g})). M _ {\mathbf {g}}\tag{11}
\]

其中，\(M_{\mathbf{f}}(\mathbf{g})\) 是如下矩阵：在每个级数元素中以 \(g_j\) 代替 \(\mathbf{X}_j\)（ \(1 \leqslant j \leqslant q\)）后得到的矩阵，这里的每个级数元素取自 \(M_{\mathbf{f}}\)；这个公式只是《代数》第 IV 章 § 5 第 8 节公式 (9) 的改写。我们以 \(M_{\mathbf{f}}(0)\) 表示 \(M_{\mathbf{f}}\) 各元素的常数项组成的矩阵；于是由 (11) 得到

\[
M _ {\mathbf {f} \circ \mathbf {g}} (0) = M _ {\mathbf {f}} (0). M _ {\mathbf {g}} (0).\tag{12}
\]

给定整数 \(n > 0\)，记

\[
\mathbf {1} _ {n} = \mathbf {X} = (\mathrm{X} _ {1}, \dots , \mathrm{X} _ {n}) \in (\mathrm{A} [ [ \mathrm{X} _ {1}, \dots , \mathrm{X} _ {n} ] ]) ^ {n},\tag{13}
\]

并把它看作只有一列的矩阵。

对于每个系统 \(\mathbf{f} = (f_1, \ldots, f_n) \in (\mathrm{A}[[\mathrm{X}_1, \ldots, \mathrm{X}_n]])^n\)，\(M_{\mathbf{f}}\) 是 \(n\) 阶方阵；我们以 \(\mathrm{J}_{\mathbf{f}}\) 或 \(\mathrm{J}_{\mathbf{f}}(\mathbf{X})\) 表示其行列式，以 \(\mathrm{J}_{\mathbf{f}}(0)\) 表示 \(\mathrm{J}_{\mathbf{f}}\) 的常数项，它等于 \(\det(M_{\mathbf{f}}(0))\)；若 \(\mathbf{g} = (g_1, \ldots, g_n)\) 是 \((\mathrm{A}[[\mathrm{X}_1, \ldots, \mathrm{X}_n]])^n\) 中的无常数项系统，则由 (11) 和 (12) 有

\[
J _ {f \circ g} = J _ {f} (g) \cdot J _ {g}\tag{14}
\]

\[
\mathrm{J} _ {\mathbf {f} \circ \mathbf {g}} (0) = \mathrm{J} _ {\mathbf {f}} (0) \mathrm{J} _ {\mathbf {g}} (0).\tag{15}
\]

命题 5. 设 A 为交换环，\(\mathbf{f} = (f_1, \ldots, f_n)\) 是由 \(n\) 个级数组成的、属于 \(\mathrm{A}[[\mathrm{X}_1, \ldots, \mathrm{X}_n]]\) 的无常数项系统。假设 \(\mathrm{J}_{\mathbf{f}}(0)\) 在 A 中可逆。则存在无常数项系统 \(\mathbf{g} = (g_1, \ldots, g_n)\)，它由 \(n\) 个级数组成，属于 \(\mathrm{A}[[\mathrm{X}_1, \ldots, \mathrm{X}_n]]\)，使得

\[
\mathbf {f} \circ \mathbf {g} = \mathbf {1} _ {n}.\tag{16}
\]

该系统唯一，且

\[
\mathbf {g} \circ \mathbf {f} = \mathbf {1} _ {n}.\tag{17}
\]

\(\mathbf{g}\) 的存在性和唯一性由《代数》第 IV 章 § 5 第 9 节命题 10 得出；该命题应用于以下 \(n\) 个形式幂级数

\[
f _ {i} (\mathrm{Y} _ {1}, \dots , \mathrm{Y} _ {n}) - \mathrm{X} _ {i} \quad (1 \leqslant i \leqslant n).
\]

由 (15) 和 (16) 得 \(\mathrm{J}_{\mathbf{f}}(0)\mathrm{J}_{\mathbf{g}}(0)=1\)，因而 \(\mathrm{J}_{\mathbf{g}}(0)\) 也可逆。于是存在系统 \(\mathbf{h}=(h_{1},\ldots,h_{n})\)，它由 \(\mathrm{A}[[\mathrm{X}_{1},\ldots,\mathrm{X}_{n}]]\) 中的 n 个无常数项级数组成，满足 \(g\circ h=1_{n}\)；由此关系和 (16)，借助 (10)，得到

\[
\mathbf {h} = \mathbf {1} _ {n} \circ \mathbf {h} = (\mathbf {f} \circ \mathbf {g}) \circ \mathbf {h} = \mathbf {f} \circ (\mathbf {g} \circ \mathbf {h}) = \mathbf {f} \circ \mathbf {1} _ {n} = \mathbf {f}.
\]

命题 5 以及公式 (10)、(15) 表明，在合成律 \((\mathbf{f},\mathbf{g})\mapsto \mathbf{f}\circ \mathbf{g}\) 下，所有系统 \(\mathbf{f} = (f_1,\dots ,f_n)\)（它们由 \(n\) 个无常数项级数组成，属于 \(\mathrm{A}[[\mathrm{X}_1,\dots ,\mathrm{X}_n]]\)，并满足 \(\mathrm{J}_{\mathbf{f}}(0)\) 在 A 中可逆）构成一个群。

